% 张三慧《大学物理学》第四版 V1 第9章 温度和气体动理论
% 依用户提供的教材扫描PDF物理页335–337逐题重述；原书习题9.1–9.29
% 与前页「思考题9.1–9.17」区分，待专家审核和出版许可。
\section*{第9章\quad 温度和气体动理论}
\subsection*{章末习题9.1--9.29}

\paragraph{习题 9.1}
定体气体温度计的测温气泡浸入水三相点槽时，封闭气体压强
为 \(6.65\times10^3\,\mathrm{Pa}\)。
（1）用此温度计测得温度 \(373.15\,\mathrm K\) 时，压强为多少？
（2）压强降至 \(2.20\times10^3\,\mathrm{Pa}\) 时，
温度是摄氏多少度？

\paragraph{习题 9.2}
在温度差不大的情况下，牛顿冷却定律为
\(-dT/dt=A(T-T_s)\)，环境温度 \(T_s\) 保持不变。
（1）由初始温度 \(T_1\) 解出
\(T(t)=T_s+(T_1-T_s)e^{-At}\)。
（2）冬日房间停暖后初温 \(25^\circ\mathrm C\)，
室外温度 \(-10^\circ\mathrm C\)，
\(40\,\mathrm{min}\) 后室温 \(20^\circ\mathrm C\)；
问室温降至 \(15^\circ\mathrm C\) 共需多久。

\paragraph{习题 9.3}
自行车车轮直径为 \(71.12\,\mathrm{cm}\)（28英寸），
轮胎截面直径为 \(3\,\mathrm{cm}\)；
打气筒长 \(30\,\mathrm{cm}\)、截面半径 \(1.5\,\mathrm{cm}\)。
打气20次后问轮胎内气体压强。
环境气温 \(-3^\circ\mathrm C\)，
打足后胎内温度为 \(7^\circ\mathrm C\)。
取环境初压 \(1\,\mathrm{atm}\)，理想气体及等量吸气模型。

\paragraph{习题 9.4}
在 \(90\,\mathrm{km}\) 高空，气压为
\(0.18\,\mathrm{Pa}\)，密度为
\(3.2\times10^{-6}\,\mathrm{kg/m^3}\)。
求此处的温度与分子数密度。
空气平均摩尔质量 \(29.0\,\mathrm{g/mol}\)。

\paragraph{习题 9.5}
大热气球容积 \(2.1\times10^4\,\mathrm{m^3}\)，
气球本身及负载总质量 \(4.5\times10^3\,\mathrm{kg}\)。
外部空气温度 \(20^\circ\mathrm C\)，
为使气球至少开始上升，球内空气最低应加热到多少度？
内外气体均视为大气压下的空气。

\paragraph{习题 9.6}
目前能够获得的极限真空压强为
\(1.00\times10^{-18}\,\mathrm{atm}\)。
温度 \(20^\circ\mathrm C\) 时，
体积 \(1\,\mathrm{cm^3}\) 中平均有多少个分子？

\paragraph{习题 9.7}
星际空间氢云的氢原子数密度
\(10^{10}\,\mathrm{m^{-3}}\)，温度可达
\(10^4\,\mathrm K\)。
估算氢云气体压强。

\paragraph{习题 9.8}
较高范围内大气温度随高度近似线性下降：
\(T(y)=T_0-\alpha y\)，其中 \(T_0\) 是地面温度。
（1）证明静力平衡大气压随高度关系
\[
 p=p_0\exp\!\left[
 \frac{Mg}{\alpha R}\ln\!\left(1-\frac{\alpha y}{T_0}\right)
 \right].
\]
（2）证明当 \(\alpha\to0\) 时退化为等温大气压公式。
（3）取 \(\alpha=0.6^\circ\mathrm C/(100\,\mathrm m)\)，
\(M=29.0\,\mathrm{g/mol}\)，
\(T_0=273\,\mathrm K\)，
\(p_0=1.00\,\mathrm{atm}\)，求珠穆朗玛峰顶的
温度及气压（峰高可取 \(8848\,\mathrm m\) 作近似）。

\paragraph{习题 9.9}
证明热平衡时，两分子相对速率的平均值
\(\overline{u}\) 与单个分子平均速率
\(\overline v\) 之间有
\(\overline u=\sqrt2\,\overline v\)。
假设两分子独立，服从同一温度的麦克斯韦分布。

\paragraph{习题 9.10}
不论速率分布函数形式如何，
证明速率方均根
\(\sqrt{\overline{v^2}}\ge\overline v\)。

\paragraph{习题 9.11}
氮分子的有效直径
\(d=3.8\times10^{-10}\,\mathrm m\)。
求标准状态下的平均自由程和连续两次碰撞的
平均时间。

\paragraph{习题 9.12}
真空管线度约 \(10^{-2}\,\mathrm m\)，
管内压强 \(1.33\times10^{-3}\,\mathrm{Pa}\)，
气体分子有效直径 \(3.0\times10^{-10}\,\mathrm m\)。
求 \(27^\circ\mathrm C\) 时单位体积的分子数、
平均自由程和平均碰撞频率。

\paragraph{习题 9.13}
在 \(160\,\mathrm{km}\) 高空，
空气密度 \(1.5\times10^{-9}\,\mathrm{kg/m^3}\)，
温度 \(500\,\mathrm K\)，分子直径
\(3.0\times10^{-10}\,\mathrm m\)。
估算空气分子的平均自由程与连续两次碰撞的
平均时间，摩尔质量近似取 \(29\,\mathrm{g/mol}\)。

\paragraph{习题 9.14}
在气体放电管内，电子热运动速率
比气体分子的平均速率大很多，
可把气体分子近似视为静止靶。
若忽略电子的有效直径，与气体分子有效直径
\(d\) 碰撞，证明电子平均自由程为
\(\lambda_e=1/(n\pi d^2)\)；
同时给出碰撞横截面积。

\paragraph{习题 9.15}
热气球充氮气 \(8.5\,\mathrm g\)，温度
\(17^\circ\mathrm C\)，
在空中以 \(65\,\mathrm{km/h}\) 的速度飞行。
求（1）氮气的平动平均动能、平均转动动能和
平均总动能；（2）球内氮气内能；
（3）球内氮气随气球整体运动的轨道动能。
把常温 \(\mathrm{N_2}\) 作为双原子理想气体，
不激发振动自由度。

\paragraph{习题 9.16}
温度 \(27^\circ\mathrm C\) 时，
\(1\,\mathrm{mol}\) 氢气、氮气和氧气各有多少内能？
各取 \(1\,\mathrm g\) 时，各有多少内能？
常温下都视为双原子理想气体。

\paragraph{习题 9.17}
容器中央的隔板分开体积相等的两半。
一边装氢气温度 \(250\,\mathrm K\)，
另一边装氧气温度 \(310\,\mathrm K\)，
两边原本压强相等。
抽去隔板后充分混合、绝热且容器刚性，
求最终平衡温度。

\paragraph{习题 9.18}
\(N\) 个粒子的速率分布为
\[
 f(v)=\begin{cases}
 av/v_0,&0\le v<v_0,\\
 a,&v_0\le v\le2v_0,\\
 0,&v>2v_0 .
 \end{cases}
\]
（1）画分布曲线并求常数 \(a\)；
（2）求速率大于和小于 \(v_0\) 的粒子数；
（3）求平均速率。

\paragraph{习题 9.19}
日冕温度 \(2.0\times10^6\,\mathrm K\)，
求其中电子的方均根速率。
星际空间温度 \(2.7\,\mathrm K\)，
主要是氢原子，求氢原子的方均根速率。
激光冷却的钠原子温度
\(2.4\times10^{-11}\,\mathrm K\)，
求其方均根速率。

\paragraph{习题 9.20}
火星质量是地球的 \(0.108\) 倍，
半径是地球的 \(0.531\) 倍。
求火星表面逃逸速率；
在表面温度 \(240\,\mathrm K\) 下，
分别求 \(\mathrm{CO_2}\) 与 \(\mathrm{H_2}\)
分子的方均根速率，解释火星大气以
\(\mathrm{CO_2}\) 为主。
木星质量是地球的318倍，
半径为地球的11.2倍；
求其逃逸速率及 \(130\,\mathrm K\)
时氢分子的方均根速率，解释其富氢大气。

\paragraph{习题 9.21}
烟粒是悬浮于空气中的小颗粒，
质量 \(1.6\times10^{-16}\,\mathrm{kg}\)。
估算 \(300\,\mathrm K\) 时烟粒的方均根速率。
若改悬浮于同温氢气中，它的方均根速率
是否不同？说明原因。

\paragraph{习题 9.22}
悬浮于 \(27^\circ\mathrm C\) 液体中的碳粒，
质量 \(6.2\times10^{-14}\,\mathrm g\)，
观测方均根速率 \(1.4\,\mathrm{cm/s}\)。
利用气体普适常量 \(R\) 与能量均分定理
求阿伏伽德罗常数的实验估计。

\paragraph{习题 9.23}
利用麦克斯韦速率分布推导平动能
\(\epsilon=\tfrac12mv^2\) 的概率分布：
\[
 f_\epsilon(\epsilon)d\epsilon=
 \frac2{\sqrt\pi}(k_BT)^{-3/2}
 \epsilon^{1/2}e^{-\epsilon/(k_BT)}d\epsilon.
\]
求最概然平动能，并与
\(\tfrac12m v_p^2\) 比较；\(v_p\) 为
最概然\emph{速率}。

\paragraph{习题 9.24}
佩兰实验中，藤黄粒子在水中的分布近似满足
\[
 \frac{RT}{N_A}\ln\frac{n_0}{n}
 =\frac43\pi a^3(\Delta-\delta)gh,
\]
\(\Delta,\delta\) 分别为藤黄、水密度，
\(a\) 是粒子半径。
（1）从玻尔兹曼高度分布推导此式。
（2）取 \(a=0.212\times10^{-6}\,\mathrm m\)，
\(\Delta-\delta=0.2067\,\mathrm{g/cm^3}\)，
\(t=20^\circ\mathrm C\)，
每相隔 \(30\times10^{-6}\,\mathrm m\) 的
四个液层粒子数依次为
\(7160,3360,1620,860\)；
通过第一与第二、第三层统计推断
阿伏伽德罗常数并说明抽样误差。

\paragraph{习题 9.25}
密闭气缸内有液态水和饱和水蒸气，
温度 \(100^\circ\mathrm C\)、压强 \(1\,\mathrm{atm}\)，
水蒸气摩尔体积
\(3.01\times10^4\,\mathrm{cm^3/mol}\)。
（1）求每 \(\mathrm{cm^3}\) 水蒸气中的分子数；
（2）求每秒撞击 \(1\,\mathrm{cm^2}\) 液面的蒸气分子数；
（3）平衡时每秒每 \(\mathrm{cm^2}\) 液面逸出的
水分子数；
（4）等温活塞将水蒸气体积缩小一半后，
蒸气压强为多少？
\textit{仍有足够液态水与液汽平衡。}

\paragraph{习题 9.26}
容器容积 \(20\,\mathrm L\)，内有
\(1.1\,\mathrm{kg}\) \(\mathrm{CO_2}\)，
温度 \(13^\circ\mathrm C\)。
用范德瓦耳斯状态方程求压强，
给定
\(a=3.64\times10^5\,\mathrm{Pa\,L^2/mol^2}\)、
\(b=0.0427\,\mathrm{L/mol}\)；
与理想气体压强比较并讨论体积排斥和吸引效应。

\paragraph{习题 9.27}
参照教材图9.18、9.21，
\(\mathrm{CO_2}\) 的范德瓦耳斯等温线
在临界点 \(K\) 有水平拐点。
（1）把范德瓦耳斯方程写为
\(p=p(V_m)\)，利用临界点
\(\partial p/\partial V_m=
 \partial^2p/\partial V_m^2=0\)
推导
\(a=3V_{m,c}^2p_c,\ b=V_{m,c}/3\)。
（2）根据教材图示临界点压强和摩尔体积
估算 \(a,b\)（图值读取精度须标记）。

\paragraph{习题 9.28}
在标准状态下，氦气的黏度
\(\eta=1.89\times10^{-5}\,\mathrm{Pa\,s}\)，
摩尔质量
\(M=0.004\,\mathrm{kg/mol}\)，
平均速率
\(\overline v=1.20\times10^3\,\mathrm{m/s}\)。
求氦原子平均自由程和有效直径，
按稀薄气体动力论模型计算。

\paragraph{习题 9.29}
热水瓶两壁间隔
\(l=0.4\,\mathrm{cm}\)，
间隙充氮气，温度 \(27^\circ\mathrm C\)，
有效分子直径
\(d=3.7\times10^{-10}\,\mathrm m\)。
当间隙压强低到多少时，
气体导热率开始明显小于通常常压值？
用平均自由程与间距相当作为稀薄气体判据。
